\section{Optimization}

The objective in Section~3 is nonconvex in the Tucker core, factors, and link
coefficients jointly, but it has a useful block structure. With the polynomial
link fixed, the loss can be differentiated through the Tucker reconstruction.
With the latent tensor fixed, the link coefficients solve a low-dimensional
ridge regression on the currently observed entries. Algorithm~\ref{alg:ntdpl}
summarizes the solver used in the experiments.

\begin{algorithm}[t]
  \caption{Masked NTD-PL optimization}
  \label{alg:ntdpl}
  \begin{algorithmic}[1]
    \Require observed tensor \(\cY\), mask \(\Omega\), rank \((R_1,\ldots,R_N)\), degree \(P\)
    \State initialize \(\cX,\{\bA^{(n)}\}\) from Tucker on the masked tensor
    \State initialize active degree \(P_{\mathrm{act}}\leftarrow 1\) and
    \(\bbeta=(0,1)\)
    \For{\(t=1,\ldots,T\)}
      \State form \(\cS=\tucker{\cX}{\bA^{(1)},\ldots,\bA^{(N)}}\)
      \If{a link update is scheduled}
        \State update active coefficients
        \(\bbeta_{0:P_{\mathrm{act}}}\) by ridge regression on
        \(\{(S_{\mathbf i},Y_{\mathbf i}):\Omega_{\mathbf i}=1\}\)
      \EndIf
      \State compute \(f_{\bbeta}(\cS)\) and \(f_{\bbeta}'(\cS)\)
      \State take first-order steps on \(\cX\) and \(\{\bA^{(n)}\}\) using
      the masked gradient
      \State normalize factor scales and absorb them into the core
      \State increase \(P_{\mathrm{act}}\) according to the continuation schedule,
      initializing the new coefficient at zero
    \EndFor
    \State \Return \(\cX,\{\bA^{(n)}\},\bbeta\)
  \end{algorithmic}
\end{algorithm}

Unlike a post-hoc calibration baseline, the Tucker tensor is not held fixed
after initialization. The link update uses the current latent tensor, and the
subsequent backbone gradient step is taken under the current nonlinear link.
This coupling is the optimization counterpart of the model distinction made
in Section~3.

\subsection{Backbone Gradient Step}

For a fixed \(\bbeta\), define
\begin{equation}
  f_{\bbeta}'(\cS)
  =
  \sum_{p=1}^{P}p\beta_p\cS^{\odot(p-1)}.
\end{equation}
The gradient of the normalized masked data term with respect to the latent
tensor \(\cS\) is
\begin{equation}
  \bT =
  \frac{1}{|\Omega|}
  \Omega \odot
  \left(f_{\bbeta}(\cS)-\cY\right)
  \odot f_{\bbeta}'(\cS).
  \label{eq:latent-gradient}
\end{equation}
Backpropagating \(\bT\) through the Tucker reconstruction gives the core
gradient
\begin{equation}
  \nabla_{\cX}\mathcal L
  =
  \bT
  \times_1 \bA^{(1)\top}
  \times_2 \cdots
  \times_N \bA^{(N)\top}
  +
  \lambda_X\cX.
\end{equation}
For the \(n\)-th factor, let
\begin{equation}
  \mathbf{Z}^{(n)}
  =
  \left[
  \cX
  \times_1 \bA^{(1)}
  \cdots
  \times_{n-1}\bA^{(n-1)}
  \times_{n+1}\bA^{(n+1)}
  \cdots
  \times_N \bA^{(N)}
  \right]_{(n)} .
\end{equation}
Then the factor gradient can be written in unfolding form as
\begin{equation}
  \nabla_{\bA^{(n)}}\mathcal L
  =
  \bT_{(n)}\mathbf{Z}^{(n)\top}
  +
  \lambda_n\bA^{(n)}.
\end{equation}
The experiments use first-order updates for these gradients. The algorithmic
claim does not depend on a particular adaptive optimizer; the important point
is that the backbone is updated jointly with the current nonlinear link.

\subsection{Closed-Form Link Update}

With \(\cS\) fixed and active degree \(Q=P_{\mathrm{act}}\), the coefficient
subproblem is ridge regression. Let \(\by_\Omega\) collect the observed values
and let \(\bPhi_{\Omega,Q}(\cS)\) have rows
\([1,s,s^2,\ldots,s^Q]\) for \(s=S_{\mathbf i}\),
\(\Omega_{\mathbf i}=1\). The subproblem corresponding to the normalized
objective in Section~3 is
\begin{equation}
  \min_{\bbeta_{0:Q}}
  \frac{1}{2|\Omega|}
  \norm{\by_\Omega-\bPhi_{\Omega,Q}(\cS)\bbeta_{0:Q}}_2^2
  +
  \frac{\lambda_\beta}{2}\norm{\bbeta_{0:Q}}_2^2.
\end{equation}
Hence
\begin{equation}
  \bbeta_{0:Q}^\star =
  \left(
  \bPhi_{\Omega,Q}^\top\bPhi_{\Omega,Q}
  +
  |\Omega|\lambda_\beta\bI
  \right)^{-1}
  \bPhi_{\Omega,Q}^\top\by_\Omega.
  \label{eq:beta-ridge-update}
\end{equation}
Equivalently, one may absorb \(|\Omega|\) into the ridge parameter and solve
the usual unnormalized normal equations. If the constant term is disabled,
the first column is removed and \(\beta_0\) is fixed at zero. In practice the
least-squares system is tiny because \(Q\le P\le6\) in our experiments.

\subsection{Continuation and Stabilization}

The polynomial family is nested in the maximum degree, so we use a continuation
schedule that starts from the Tucker-compatible degree-one model and activates
one additional degree at a time:
\[
  t_k = \operatorname{round}\left(\frac{kT}{P}\right),
  \qquad k=1,\ldots,P-1.
\]
When degree \(p+1\) is activated, the previous coefficients and Tucker factors
are kept and the new coefficient \(\beta_{p+1}\) is initialized at zero. The
model therefore expands from a linear low-rank fit instead of starting from a
high-degree polynomial link immediately.

This staging is useful because high powers of \(\cS\) are sensitive to scale
and can dominate early gradients. We also normalize factor columns during
training and absorb their scales into the core; this is a gauge fixing of the
Tucker parameterization and does not change \(\cS\). For numerical stability
in the ridge update, the implementation may solve the Vandermonde system on a
centered and scaled latent variable \((S-\mu)/\sigma\), then convert the
coefficients back to the original power basis of \(S\). This changes only the
conditioning of the coefficient solve, not the model being optimized.

\subsection{Cost}

The dominant cost is forming the Tucker reconstruction and its reverse
contractions on the observed tensor. For dense tensors this is comparable to
standard Tucker optimization at the same rank, plus \(O(P|\Omega|)\) work to
evaluate the scalar polynomial and its derivative. The link update forms a
Vandermonde design or the corresponding moment equations and solves a
\((P+1)\times(P+1)\) linear system, with direct cost
\(O(P^2|\Omega|+P^3)\). Since \(P\) is small, this cost is negligible compared
with the Tucker contractions. Thus the polynomial link changes the observation
model without introducing a tensor-valued high-order parameter block.
